% ---
% id: statistics/singleton/thm:singleton-completeness
% kind: thm
% title: "Completeness of a singleton acceptance set"
% summary: >-
%   Under i.i.d. honest rounds a singleton acceptance set accepts with its
%   point's multinomial probability. That is 0, with no completeness budget
%   below one, if n times a coordinate is no integer or the point leaves the
%   law's support; else the least budget is one minus it, zero only for one
%   point mass.
% prominence: main
% status: implemented
% uses: [statistics/def:freq, statistics/def:matched-alphabet]
% mentions: [req:singleton-budget-record]
% assumes: [A-HON, A-ALPHABET]
% refs: []
% provenance: newly-derived
% implements: certiqs_key.statistics.singleton:singleton_completeness
% lean:
% ---

A test that accepts exactly one \tref{frequency vector} is far more
restrictive than a box around the honest law. Under independent honest
rounds with one law its accept probability is a single multinomial mass, so
its completeness can be computed exactly; an \tref{instance} that declares such
a set records the least budget of (ii), or that none exists
(Requirement~\ref{req:singleton-budget-record}).

\begin{theorem}\label{thm:singleton-completeness}
Assume A-HON, let $\Chatset$ and $n$ be a
\tref{matched alphabet and round count}
(Definition~\ref{def:matched-alphabet}), and let
$\acSet=\{\mathbf{q}^{\ast}\}$ be a singleton geometry with
$\mathbf{q}^{\ast}\in\simplex{\Chatset}$. Call a completeness
budget $\ecomp$ \emph{admissible for this geometry} when $\ecomp\in[0,1)$ and
$\Pr[\operatorname{freq}_{\Chat_1^n}\notin\acSet]\le\ecomp$.
\begin{enumerate}[label=(\roman*),leftmargin=2.2em]
\item \emph{(exact accept probability)} With the convention $0^{0}=1$,
\begin{equation}\label{eq:singleton-acc}
  \Pr\big[\operatorname{freq}_{\Chat_1^n}\in\acSet\big]
  =\begin{dcases}
     \frac{n!}{\prod_{\hat c\in\Chatset}\big(n\,q^{\ast}(\hat c)\big)!}
     \prod_{\hat c\in\Chatset}\honest(\hat c)^{\,n\,q^{\ast}(\hat c)},
     &\text{if }n\,q^{\ast}(\hat c)\in\mathbb{Z}_{\ge0}\ \forall\hat c,\\
     0,&\text{otherwise;}
   \end{dcases}
\end{equation}
\item \emph{(the least admissible budget, and when there is none)} the accept
probability \eqref{eq:singleton-acc} is $0$ if and only if either
$n\,q^{\ast}(\hat c)\notin\mathbb{Z}_{\ge0}$ for some $\hat c$, or
$\mathbf{q}^{\ast}\not\ll\honest$, that is, some $\hat c$ has
$q^{\ast}(\hat c)>0$ and $\honest(\hat c)=0$. In that case \emph{no}
$\ecomp\in[0,1)$ is admissible for this geometry, and no completeness
statement whatsoever can be made about it; a budget of $1$ lies outside
$[0,1)$. Otherwise the accept probability is strictly positive and the least
admissible completeness budget is
\begin{equation}\label{eq:singleton-budget}
  \ecomp=1-\Pr\big[\operatorname{freq}_{\Chat_1^n}\in\acSet\big]\ \in[0,1),
\end{equation}
which does lie in $[0,1)$;
\item \emph{(when the budget is zero)} the budget \eqref{eq:singleton-budget}
equals $0$ if and only if $\honest$ is a point mass $\delta_{\hat c_0}$ for
some $\hat c_0\in\Chatset$ and $\mathbf{q}^{\ast}=\delta_{\hat c_0}$; in
particular it equals $0$ whenever $|\Chatset|=1$, and it is strictly positive
whenever $\honest$ gives positive mass to two distinct outcomes.
\end{enumerate}
\end{theorem}

\begin{proof}
\emph{(i).} By Definition~\ref{def:freq}, $\operatorname{freq}_{\Chat_1^n}=
\mathbf{q}^{\ast}$ holds exactly when $\Ncount{\hat c}=n\,q^{\ast}(\hat c)$ for
every $\hat c$, which is impossible unless every $n\,q^{\ast}(\hat c)$ is a
non-negative integer; that gives the second branch. In the first branch write
$n_{\hat c}\coloneqq n\,q^{\ast}(\hat c)$, so $\sum_{\hat c}n_{\hat c}=n$. Under
A-HON each string $\hat c_1^n$ has probability
$\prod_{\hat c}\honest(\hat c)^{\,n_{\hat c}}$ whenever its count vector is
$(n_{\hat c})_{\hat c}$, and the number of such strings is the multinomial
coefficient $n!/\prod_{\hat c}n_{\hat c}!$; summing over them gives
\eqref{eq:singleton-acc}.

\emph{(ii).} If some $n\,q^{\ast}(\hat c)$ is not a non-negative integer the
second branch of \eqref{eq:singleton-acc} gives probability $0$. Otherwise the
first branch applies; its multinomial coefficient is a positive integer, and
under the convention $0^{0}=1$ each factor $\honest(\hat c)^{\,n_{\hat c}}$
with $n_{\hat c}=0$ equals $1$. The product therefore vanishes if and only if
some $\hat c$ has $n_{\hat c}\ge1$ and $\honest(\hat c)=0$; and $n_{\hat c}\ge1$
is equivalent to $q^{\ast}(\hat c)>0$, so this says precisely that
$\mathbf{q}^{\ast}\not\ll\honest$. That establishes the stated
characterization of a zero accept probability.

By the definition in the statement, a budget $\ecomp$ is admissible for this
geometry exactly when it satisfies both the completeness requirement
$\Pr[\operatorname{freq}_{\Chat_1^n}\notin\acSet]\le\ecomp$, that is,
$\ecomp\ge1-\Pr[\operatorname{freq}_{\Chat_1^n}\in\acSet]$, and the domain
$\ecomp\in[0,1)$ of the completeness budget. The admissible budgets
are therefore the interval
$[\,1-\Pr[\operatorname{freq}_{\Chat_1^n}\in\acSet],\,1)$, which is non-empty
if and only if $\Pr[\operatorname{freq}_{\Chat_1^n}\in\acSet]>0$, and which
then has least element \eqref{eq:singleton-budget}. When the accept
probability is $0$ the interval is empty: there is no admissible budget, and
in particular $\ecomp=1$ is not one.

\emph{(iii).} The budget \eqref{eq:singleton-budget} is $0$ exactly when the
accept probability is $1$. That probability equals $1$ if and only if
$\operatorname{freq}_{\Chat_1^n}=\mathbf{q}^{\ast}$ almost surely. Suppose
$\honest$ gives positive mass to two distinct outcomes $\hat c\neq\hat c'$.
The string $(\hat c,\hat c,\dots,\hat c)$ and the string
$(\hat c',\hat c,\dots,\hat c)$ both have positive probability under A-HON, and
their frequency vectors differ (they differ by $1/n$ in coordinate $\hat c$),
so $\operatorname{freq}_{\Chat_1^n}$ is not almost surely constant and the
probability is strictly less than $1$. Hence $\honest$ must be a point mass
$\delta_{\hat c_0}$; then $\operatorname{freq}_{\Chat_1^n}=\delta_{\hat c_0}$
almost surely and the probability is $1$ exactly when
$\mathbf{q}^{\ast}=\delta_{\hat c_0}$. If $|\Chatset|=1$ the only distribution
on $\Chatset$ is that point mass, and $\mathbf{q}^{\ast}$ equals it.
\end{proof}
